\documentclass[11pt,a4paper]{article}
\usepackage[utf8]{inputenc}
\usepackage[T1]{fontenc}
\usepackage[margin=0.75in]{geometry}
\usepackage{enumitem}
\usepackage{titlesec}
\usepackage{xcolor}
\usepackage{hyperref}

\definecolor{primarycolor}{RGB}{33,37,41}
\definecolor{secondarycolor}{RGB}{0,102,204}
\definecolor{accentcolor}{RGB}{108,117,125}
\hypersetup{colorlinks=true,urlcolor=secondarycolor,linkcolor=secondarycolor}
\titleformat{\section}{\large\bfseries\color{primarycolor}}{}{0em}{}
\titlespacing*{\section}{0pt}{1.1em}{0.45em}
\setlist[itemize]{leftmargin=*,itemsep=0.25em,topsep=0.25em}
\pagestyle{empty}

\begin{document}

\begin{center}
{\Huge\bfseries\color{primarycolor} MARTA KOWAL-BRANDT}\\[0.3em]
{\large QA Automation Engineer | CI/CD \& Test Infrastructure}\\[0.25em]
Wroclaw, Poland | marta.kowalbrandt@example.com | +48 555 303 303 | gitlab.com/mkb-example
\end{center}

\section*{PROFESSIONAL SUMMARY}
QA automation engineer with 7 years of experience, including the last two focused primarily on CI/CD pipelines, test environments and build tooling. Owns GitLab CI pipelines supporting three product teams and has improved pipeline speed, test reliability and the availability of ephemeral test environments. Brings hands-on experience with Docker, Helm, Kubernetes, Linux, Python and Bash, and is seeking to apply this infrastructure experience in a DevOps or platform engineering role.

\section*{SKILLS}
CI/CD and build tooling: GitLab CI, Jenkins, pipeline design and maintenance, artifact publishing\\
Infrastructure and environments: Docker, Docker Compose, Helm, Kubernetes (basic), Linux\\
Programming and scripting: Python, Bash, Java (test code), SQL\\
Testing and observability: pytest, Playwright, Selenium, Postman, Grafana dashboards for test metrics\\
Databases and tools: PostgreSQL, Git

\section*{PROFESSIONAL EXPERIENCE}
\textbf{QA Automation Engineer (Test Infrastructure) | Odrava Health Software, Wroclaw | 2021--Present}
\begin{itemize}
\item Own GitLab CI pipelines for three product teams (approximately 25 developers), covering build, unit, integration and end-to-end test stages, as well as artifact publishing.
\item Reduced average pipeline duration from 38 to 14 minutes by parallelising test jobs, caching dependencies and moving end-to-end tests to on-demand ephemeral environments.
\item Built ephemeral test environments for each merge request using Docker Compose, later using Helm charts on the company’s shared Kubernetes cluster.
\item Maintain a pytest and Playwright end-to-end suite of approximately 900 tests; reduced the flaky-test rate from 7\% to under 1\%.
\item Developed Bash and Python tooling to seed test data into PostgreSQL.
\end{itemize}

\textbf{QA Engineer | Brightmill Software House, Wroclaw | 2019--2021}
\begin{itemize}
\item Performed manual and automated testing of web and mobile applications for client projects.
\item Introduced a Selenium and Java automated regression suite for the largest client project.
\item Configured Jenkins jobs for nightly regression runs.
\end{itemize}

\textbf{Customer Support Specialist (2nd line) | Brightmill Software House | 2018--2019}
\begin{itemize}
\item Investigated customer issues using logs and databases, and reproduced bugs for developers.
\end{itemize}

\section*{EDUCATION}
BEng, Computer Engineering | Silesian Lowlands University of Technology | 2018

\section*{LANGUAGES}
Polish (native), English (C1), German (A2)

\end{document}